% Notation macros.
%
% The names follow the reference documents, NOT transformer convention. In particular
% `d` is the size of the label set of a word variable, not a hidden dimension, and `h`
% counts channels (head variables per word), not attention heads in the usual sense.
% Keep this file in sync with src/config.py, which uses the same names.
%
% Every macro below must correspond to something that exists in the factor graph — the
% modelling constraint of the project is that no parameter names a non-factor, and the
% notation should not be able to express one.

% --- index sets and sizes ---
\newcommand{\Vocab}{\mathcal{V}}                 % vocabulary
\newcommand{\nlab}{d}                            % number of labels of a word variable
\newcommand{\nchan}{h}                           % number of channels (head variables per word)
\newcommand{\krank}{r}                           % Kruskal rank of the arc-score decomposition
\newcommand{\nglob}{m}                           % number of global variables (Appendix B.3)
\newcommand{\clip}{\gamma}                       % distance clip threshold
\newcommand{\niter}{T}                           % MFVI iterations / loop depth
\newcommand{\seqlen}{n}                          % sequence length
\newcommand{\Dep}[1]{\mathcal{D}_{#1}}           % dependency candidate set at position #1

% --- random variables ---
\newcommand{\Wv}[1]{W_{#1}}                      % word variable (LATENT, not observed)
\newcommand{\Zv}[1]{Z_{#1}}                      % label variable
\newcommand{\Hv}[2]{H_{#1}^{(#2)}}               % head variable at position #1, channel #2
\newcommand{\Gv}[1]{G_{#1}}                      % global variable

% --- factors and parameters (each names a factor; see CLAUDE.md constraint 1) ---
\newcommand{\Sfac}{S}                            % word--label factor, |V| x d; mediates BOTH directions
\newcommand{\Tfac}[1]{T^{(#1)}}                  % arc score of channel #1, T = U V^\top
\newcommand{\Ufac}[1]{U^{(#1)}}
\newcommand{\Vfac}[1]{V^{(#1)}}
\newcommand{\bfac}{b}                            % word unary; equals the LM-head bias
\newcommand{\rfac}[1]{r^{(#1)}}                  % root key of channel #1
\newcommand{\Bfac}[1]{B^{(#1)}}                  % contracted arc--label score entering the readout

% --- variational quantities ---
\newcommand{\Q}{Q}                               % variational posterior
\newcommand{\QW}[1]{\Q_{\Wv{#1}}}
\newcommand{\QZ}[1]{\Q_{\Zv{#1}}}
\newcommand{\QH}[2]{\Q_{\Hv{#1}{#2}}}
\newcommand{\Free}{\mathcal{F}}                  % mean-field free energy
\newcommand{\lamZ}{\lambda_Z}                    % entropic Frank--Wolfe message weights
\newcommand{\lamH}{\lambda_H}
\newcommand{\lamW}{\lambda_W}
\newcommand{\lamG}{\lambda_G}

% --- readout ---
\newcommand{\readout}[1]{\mu_{#1}}               % \mu_t(a): unnormalised readout score of word a
\DeclareMathOperator*{\LSE}{LSE}                 % log-sum-exp
\DeclareMathOperator*{\logcumsumexp}{logcumsumexp}

% --- generic ---
\newcommand{\R}{\mathbb{R}}
\newcommand{\E}{\mathbb{E}}
\newcommand{\Prob}{\mathbb{P}}
\newcommand{\KL}[2]{\mathrm{KL}\!\left(#1 \,\|\, #2\right)}
\newcommand{\softmax}{\operatorname{softmax}}
\newcommand{\prefix}[1]{{<}#1}                   % "the prefix before position #1"
\newcommand{\stopgrad}[1]{\left\lfloor #1 \right\rfloor}  % frozen condition: no gradient path

% --- added while writing Section 2 ---
\newcommand{\qbar}[1]{\bar{q}_{#1}}              % frozen filtering marginal of position #1
\newcommand{\Bglob}{B'}                          % global-variable factor matrix (App. B.3.3)
\newcommand{\ncomp}{K}                           % number of label components (factored labels)
\newcommand{\ind}[1]{\mathbf{1}\!\left[#1\right]}  % indicator; \mathbb{1} is not in the ACL fonts
\newcommand{\Hvar}[1]{H_{#1}}                    % the head-variable block at position #1

% --- added while writing Appendix C (proof of the trilemma) ---
% These name the objects of Part I of the technical note: the proof lives on a product of
% simplices, one per node of the factor graph, and everything turns on the tangent spaces of
% those simplices. \Vocab is already \mathcal{V}, so the node set cannot reuse that letter.
\newcommand{\Nodes}{\mathcal{N}}                 % node set: one node per label / head variable
\newcommand{\Xsp}{\mathcal{X}}                   % product of the node simplices
\newcommand{\Xint}{\mathcal{X}^{\circ}}          % its relative interior; every iterate lives here
\newcommand{\Tang}[1]{\mathcal{T}_{#1}}          % tangent space of node #1's simplex
\newcommand{\Proj}[1]{\Pi_{#1}}                  % centring projector onto \Tang{#1}

% --- added while writing Appendix B (additional variants) ---
% Both name objects that already exist in the graph; neither introduces a new symbol for a
% quantity that already had one.
\newcommand{\QG}[1]{\Q_{\Gv{#1}}}                % posterior of the global head at position #1
\newcommand{\nlabc}{d'}                          % width of ONE label component under factoring
                                                 % (App. B.4): \nlab = \ncomp \nlabc, i.e. \nlab
                                                 % stays the TOTAL label width, written D in the
                                                 % source document.

% --- added while writing Appendix A (derivations) ---
% Three names the main text uses in words but never abbreviates. \dmodel is the transformer
% hidden width, written d_h in the source document's cost table; it is NOT \nlab, and the whole
% point of Section 3.3 is that the comparison unit is \nlab <-> \dmodel. \Znorm is the local
% normaliser of a slot conditional and the per-component normaliser of the mixture readout.
% \Ent is the entropy: \Hv is already the head variable, so H cannot be reused unqualified.
\newcommand{\dmodel}{d_{\text{model}}}           % transformer hidden width (source: d_h)
\newcommand{\Znorm}{\mathcal{Z}}                 % normalising constant
\newcommand{\Ent}{\mathrm{H}}                    % entropy of a distribution
